\subsection{The full exceptional boundary}
\label{sec:full-boundary}

The log model must remember every point over the singular section.  This
section proves that relative nefness forces that property in the
non-log-canonical case.  It also identifies precisely the boundary entering
surface adjunction.  The existence of the model is an input to this section.

Let $C$ be a smooth connected complex curve with distinguished point $0$,
let $\sigma:X\to C$ be a flat family of normal Gorenstein surfaces, and let
$\Gamma\subset X$ be a section.  Assume that $X$ is normal and Gorenstein.
Consider a projective birational morphism
\[
 h:Y\longrightarrow X,\qquad q=\sigma\circ h,
\]
where $Y$ is a normal $\mathbb Q$-factorial threefold and $h$ is an
isomorphism over $X\setminus\Gamma$.  Write $H$ for the reduced sum of the
$h$-exceptional prime divisors dominating $C$.  Assume that all vertical
$h$-exceptional prime divisors lie over $0$, that every fiber $Y_c$ is a
reduced Cartier divisor, and that
\[
 Y_0=S_0+\sum_a S_a,
\]
where $S_0$ maps birationally onto $X_0$ and every $S_a$ maps to the point
$\Gamma\cap X_0$.  We assume that $Y_c$ is irreducible for $c\ne0$, that
$(Y,H+Y_c)$ is dlt for all $c$ under consideration, and that
\[
 L=K_Y+H
\]
is $h$-nef.  We use compatible canonical divisors and put
\[
 A=L-h^*K_X,\qquad D=Y_0+H.
\]
The assumptions on the fibers may equally be imposed after shrinking $C$.
For a component $S$ of $Y_c$, let $B_S$ be the reduced union of the curves
in which $S$ meets the other components of $H+Y_c$.

\begin{lemma}[Surface adjunction]
\label{lem:surface-boundary}
Every component $S$ of $Y_c$ is normal.  Its different is exactly $B_S$:
\[
 (K_Y+H+Y_c)|_S\sim_{\mathbb Q}K_S+B_S.
\]
In particular $(S,B_S)$ is dlt and
\[
 L|_S\sim_{\mathbb Q}K_S+B_S.
\]
The latter divisor is nef on every curve contracted by $h|_S$.  For an
extra vertical component $S_a$, it is nef on the projective surface $S_a$.
\end{lemma}

\begin{proof}
Dlt adjunction gives normality of $S$ and a dlt pair
$(S,\operatorname{Diff}_S(H+Y_c-S))$; see
\cite[Lemma~10.2]{Fujino2011}.  We check the coefficients of the different
at every prime curve of $S$.

If a prime curve lies in another component of $H+Y_c$, its generic point
is an lc stratum.  The dlt pair is simple normal crossing there, so ordinary
adjunction gives coefficient one.  Here we use the description of dlt lc
centers as boundary strata and the defining simple-normal-crossing property
at their generic points; see
\cite[Definition~11.5 and (11.10.4)]{Kollar2023Families}.

At the generic point $\eta$ of any other prime curve of $S$, no other
component of $H+Y_c$ is present.  Thus $S$ equals the Cartier divisor
$Y_c$ near $\eta$.  The local ring $\mathcal O_{S,\eta}$ is a regular
one-dimensional local ring, since $S$ is normal.  If $t$ is a local
equation of the fiber, the two-dimensional local ring
$R=\mathcal O_{Y,\eta}$ satisfies
\[
 R/(t)=\mathcal O_{S,\eta}.
\]
Consequently its maximal ideal is generated by $t$ and one lift of a
generator of the maximal ideal of $R/(t)$.  The embedding dimension of $R$
is at most two, and $R$ is regular.  Ordinary adjunction therefore gives
coefficient zero at this curve.  This proves equality of the different
and $B_S$ as divisors; no other prime-curve contribution remains.

The line bundle of $Y_c=q^*(c)$ restricts trivially to $S$, giving the
formula with $L$.  Its degree on a curve contracted by $h|_S$ is the degree
of $L$ on that curve and is nonnegative.  Finally, $S_a$ is projective
because it is a closed subvariety of the projective fiber of $h$ over
$\Gamma\cap X_0$, and every curve on $S_a$ is contracted by $h$.
\end{proof}

\begin{proposition}[Support over a non-log-canonical section]
\label{prop:full-exceptional-boundary}
Assume in addition that the surface $X_c$ is not log canonical at
$\Gamma\cap X_c$ for general $c\in C$.  Then $A\le0$, some horizontal coefficient of $A$ is
strictly negative, and
\[
 h^{-1}(\Gamma)=\operatorname{Supp}(-A)
              =\operatorname{Supp}(H)\cup\bigcup_a S_a
\]
as sets.  In particular $B_{S_0}$ is the full reduced exceptional
preimage of $\Gamma\cap X_0$ on $S_0$, and $B_{Y_c}$ has the same property
on every $Y_c$ with $c\ne0$.  Thus
\[
 S_0\setminus\operatorname{Supp}(B_{S_0})
       \simeq X_0\setminus\Gamma,
 \qquad
 Y_c\setminus\operatorname{Supp}(B_{Y_c})
       \simeq X_c\setminus\Gamma\quad(c\ne0).
\]
\end{proposition}

\begin{proof}
The divisor $A$ is exceptional, since $h$ is an isomorphism away from
$\Gamma$, and is $h$-nef, since a pullback from $X$ has degree zero on
contracted curves.  Apply the negativity lemma to $-A$: its pushforward
is zero, so $-A$ is effective
\cite[Lemma~4.16(1)]{Fujino2011}.

Suppose every horizontal coefficient of $A$ were zero.  The support of
$A$ would then lie over $0$.  For general $c\ne0$, its restriction to
$Y_c$ would be zero.  By Lemma~\ref{lem:surface-boundary} and Cartier-fiber
adjunction on $X$, with compatible canonical divisors, this gives a
crepant comparison
\[
 K_{Y_c}+B_{Y_c}=h_c^*K_{X_c},\qquad h_c=h|_{Y_c}.
\]
The pair on the left is log canonical.  Computing discrepancies on a
common resolution now shows that $X_c$ is log canonical: the discrepancy
of a divisor exceptional over $Y_c$ is unchanged, and a boundary component
on $Y_c$ has discrepancy $-1$ over $X_c$.  This contradicts the hypothesis.
Thus at least one horizontal exceptional prime has negative coefficient.

Such a prime maps onto $\Gamma$: its image is closed by properness, is
contained in $\Gamma$, and dominates $C$.  Therefore every fiber of $h$
over $\Gamma$ meets $\operatorname{Supp}(-A)$.  The support statement in
the negativity lemma says that each fiber of a proper birational morphism
between normal varieties either lies in the support of an effective
anti-nef divisor or is disjoint from it
\cite[Lemma~4.16(2)]{Fujino2011}.  Applying it to $-A$ proves
$h^{-1}(\Gamma)\subset\operatorname{Supp}(-A)$; the reverse inclusion
holds because $A$ is exceptional over $\Gamma$.

Every prime in $\operatorname{Supp}(-A)$ is either a component of $H$ or
one of the $S_a$, by the assumptions on the exceptional primes.
Conversely these divisors map into $\Gamma$ and hence lie in the support.
Intersect this equality with $S_0$.  None of the listed divisors equals
$S_0$, and their intersections are precisely the boundary curves defining
$B_{S_0}$; dlt strata have the expected codimension.  Hence there are no
additional exceptional points or curves on $S_0$ outside this boundary.
The same argument on $Y_c$, $c\ne0$, uses only $H$.  The stated
isomorphisms follow from the assumed isomorphism away from $\Gamma$.
\end{proof}

\begin{remark}
Lemma~\ref{lem:surface-boundary} identifies a divisor on a surface.  It
does not assert that further adjunction to a boundary curve has zero
fractional different, nor that the threefold is smooth at every closed
boundary point.  Those are distinct local questions.
\end{remark}
